\chapter*{Preface}

This book aims to help the reader truly understand Quantum Computing from the
ground up, and to see clearly where it departs from classical computing.


The book grew out of my own self-learning journey into the subject. Like many
students, I began with \textit{Nielsen \& Chuang}, rightly regarded as the bible
of the field. These notes are what I wrote for myself while working through that
material, rebuilding each concept from the basics. They also draw on lectures,
research papers, books, writings by eminent professors in the field, and many
other resources I came across along the way.

The central goal of this book is to make Quantum Computing as easy to understand
as possible, without compromising on mathematical rigor. In contrast to standard
texts such as \textit{Nielsen \& Chuang} and \textit{Rieffel \& Polak}, which
serve as comprehensive references for the field, this book builds understanding
step by step for an undergraduate level. The chapters start from the basics and
build up to the major quantum algorithms, with the aim of giving the reader
strong fundamentals.

Readers are recommended to begin with the \textbf{Appendix}, which introduces the
Linear Algebra and Probability used throughout the text. For readers unfamiliar with
Linear Algebra, it provides the mathematical foundation needed for the later
chapters; for readers already familiar with it, it serves to introduce the
notation used throughout the book.

Alongside the theory, the book includes \textbf{code snippets} that show how the
ideas are implemented in practice. Companion Jupyter notebooks for many chapters
are available at
\url{https://github.com/Nihar-Shah2001/Quantum_Computing_Notes}. Each chapter
ends with the references relevant to it, and a complete list of references is
provided at the end of the book.

These notes are a continuous work in progress. I keep revising and improving the
existing chapters and adding new ones, so the book will continue to grow over
time.

I hope this text helps you build a strong and rigorous understanding of Quantum
Computing.

\begin{flushright}
\textit{--- Nihar Shah}
\end{flushright}
